\documentclass[10pt,a4paper]{amsart}
\usepackage[margin=19mm]{geometry}
\usepackage{amsmath,amssymb,mathtools}
\usepackage[scr=rsfs,cal=euler]{mathalfa}
\usepackage{xfrac}
\usepackage[unicode,hidelinks,bookmarks=false]{hyperref}
\newcommand{\ie}{i.e.\ }
\newcommand{\cf}{cf.\ }
\newcommand{\resp}{respectively\ }
\newcommand{\Subsection}{\subsection}
\providecommand{\nobreakdash}{\nobreak\hbox{-}}
\setlength{\emergencystretch}{2em}
\multlinegap0pt
\usepackage{bm}
\usepackage{bbm}
\usepackage{dsfont}
\usepackage{xspace}
\usepackage{cancel}
\usepackage{mathtools}
\usepackage{amsthm}
\makeatletter
\@ifundefined{theorem}{\newtheorem{theorem}{Theorem}}{}
\makeatother
\def\E{{\mathbb {E}}}
\newcommand{\fm}{\ensuremath{\mathfrak{m}}}
\newcommand{\fv}{\ensuremath{\mathfrak{v}}}
\title[On Card guessing after a single shelf shuffle]{On Card guessing after a single shelf shuffle}
\author{Markus Kuba}
\date{Source: arXiv:2602.12928v2 (CC BY 4.0)}
\begin{document}
\maketitle

\noindent\emph{Source and licence.} On Card guessing after a single shelf shuffle. Version arXiv:2602.12928v2 (CC BY 4.0), distributed under the Creative Commons Attribution 4.0 International licence, as recorded in the arXiv licence metadata for this exact version and shown on the abstract page https://arxiv.org/abs/2602.12928v2. Full rights notice: licenses/arxiv-2602.12928-CC-BY-4.0.txt.\par\smallskip

\noindent\textit{Editorial scope.} Counted unit: the Theorem whose source label is \texttt{None} in the article (source0.tex lines 424--459, source label \texttt{None}), delivered together with the article's own complete original proof of it (source0.tex lines 461--484, one \texttt{proof} environment of 24 lines). No mathematics is written, repaired or re-derived: statement and proof are the article's own text line for line. Required context retained uncounted: the:main (lines 406--413). Every definition, display and label of those lines is retained unchanged. Omitted from this package and not redistributed: 1--405, 414--423, 460--460, 485--816. Nothing inside a retained excerpt is omitted or altered: every retained body is delivered exactly as the article's lines are written, so no omission receipt of any kind accompanies this package. Citations: the retained excerpts carry no citation command at all, so no bibliography entry of the article is transcribed and no reference is redistributed. Reference adaptation: none was needed and none was made; every reference inside the delivered unit resolves inside it, so no unresolved reference or citation is produced. Printed slips kept literal: no printed word, symbol, hypothesis or number of the counted statement or of its proof was altered. Layout and class: the article's own class is \texttt{amsart} with packages including graphicx, pgfplots, tikz, caption, amsmath, amssymb, amsthm, hyperref, xcolor, times, enumitem, verbatim, listings, fancyvrb; the wrapper is the lane's pinned \texttt{amsart} editorial wrapper at 10pt on A4, which restates the theorem-like environments the retained text needs and adds the article's own macro definitions that the retained text needs (\texttt{\textbackslash E}, \texttt{\textbackslash fm}, \texttt{\textbackslash fv}), with extra wrapper packages (bm, bbm, dsfont, xspace, cancel, mathtools). The delivered numbering is the unit's own. Assets: no retained span contains a figure environment, a graphics-inclusion command, a PDF-page inclusion, a table or a TikZ picture; no image file is copied into this package, the package declares no asset dependency and no image file is redistributed with it. Licence: version arXiv:2602.12928v2 of this article is distributed under the Creative Commons Attribution 4.0 International licence, as recorded in the arXiv licence metadata for this exact version and shown on the abstract page https://arxiv.org/abs/2602.12928v2. A full rights notice is delivered at licenses/arxiv-2602.12928-CC-BY-4.0.txt. Characters: every retained excerpt is pure ASCII, so the delivered engine, XeLaTeX with its default Latin Modern text font, reports no missing character.
\begin{theorem}
\label{the:main}
The random variable $X=X_n$ of correctly guessed cards, starting with a deck of $n$ cards, after a one-time shelf shuffle satisfies
a central limit theorem
\[
\frac{X_n-\frac34 n}{\sqrt{n}/4}\to \mathcal{N}(0,1). 
\]
\end{theorem}

\begin{theorem}[Hwang's quasi-powers Theorem]
Let the $X_n$ be non-negative discrete random variables, supported by $\mathbb{Z}_{\ge 0}$, with probability generating functions $p_n(u)$. Assume that,
uniformly in a fixed complex neighbourhood of $v = 1$, for sequences
$\beta_n, \kappa_n \to +\infty$, one has
\begin{equation*}
\E(v^{X_n})=  A(v)\, B(v)^{\beta_n} \left( 1 + O\!\left( \frac{1}{\kappa_n} \right) \right),
\end{equation*}
where $A(u)$, $B(u)$ are analytic at $u = 1$ and $A(1) = B(1) = 1$.
Assume finally that $B(u)$ satisfies the so-called variability condition,
\[
\fv(B(v))=B''(1) + B'(1) - B'(1)^2 \neq 0.
\]
Under these conditions, the mean and variance of $X_n$ satisfy
\begin{equation*}
\mu_n = \mathbb{E}(X_n)
= \beta_n \fm(B(v)) + \fm(A(v)) + O\!\left( \kappa_n^{-1} \right),
\end{equation*}
\begin{equation*}
\sigma_n^2 = \mathbb{V}(X_n)
= \beta_n \fv(B(v)) + \fv(A(v)) + O\!\left( \kappa_n^{-1} \right).
\end{equation*}
The distribution of $X_n$ is, after standardization, asymptotically Gaussian, and the
speed of convergence to the Gaussian limit is $\mathcal{O}\left( \kappa_n^{-1} + \beta_n^{-1/2} \right)$:
\begin{equation*}
\mathbb{P}\left\{
\frac{X_n - \mathbb{E}(X_n)}{\sqrt{\mathbb{V}(X_n)}} \le x
\right\}
=\Phi(x)+\mathcal{O}\left(\frac{1}{\kappa_n}+\frac{1}{\sqrt{\beta_n}}\right),
\end{equation*}
where $\Phi(x)$ is the distribution function of a standard normal,
\[
\Phi(x)=\frac{1}{\sqrt{2\pi}}
\int_{-\infty}^{x}
e^{-v^2/2}\, dv.
\]
\end{theorem}

\begin{proof}[Proof of Theorem~\ref{the:main}]
We analyze the singular structure of $S(z,v)$. The denominator 
\[
4-4vz+(v^2-v)z
\]
has two roots $\rho_1(v)$ and $\rho_2(v)$
\[
\rho_1(v)=\frac{2(v-\sqrt{v})}{v(v-1)},
\quad \rho_2(v)=\frac{2(v+\sqrt{v})}{v(v-1)},
\]
with dominant singularity $\rho(v)=\rho_1(v)$, satisfying 
$\lim_{v\to 1} \rho(v)=1$. The quasi-power theorem applies because the dominant singularity $\rho(v)$ is simple and analytic near $v=1$,.
and $S(z,v)$ is analytic except at $z=\rho(v)$. We expand $S(z,v)$ around $z=\rho(v)$, for $v$ close to $1$, 
obtaining the singular expansion
\[
S(z,v)\sim \frac{A(v)}{1 - \frac{z}{\rho(v)}}, \quad A(v) \neq 0.
\]
By standard contour integration we get the desired expansion of $s_n(v)=\E(v^{X_n})$:
\[
\E(v^{X_n})=[z^n]S(z,v)\sim \frac{A(v)}{\rho(v)^n}.
\]
Application of the Quasi-Power theorem then leads to the stated result; one may also check that the asymptotics $\beta_n\fm(B(v))=3n/4$ and
$\beta_n\fv(B(v))=n/16$.
\end{proof}

\end{document}
